\section{Introduction}
\label{Intro}

Aircraft spend the majority of their flight at cruising altitudes (32,000--42,000~ft), where the air is thin enough to significantly reduce drag, whilst still dense enough to maintain sufficient oxygen flow to the engines. In atmospheric science terminology, this region is known as the Upper Troposphere Lower Stratosphere (UTLS), characterised by its low temperatures, low concentrations of atmospheric trace species and higher radiative efficiencies compared to ground level. Aviation is unique in that it is the only significant source of anthropogenic emissions in the UTLS.

\subsection{Aviation Non-\texorpdfstring{\ce{CO_2}}{CO2} Climate Impact}
The two most significant contributions to aviation's high-altitude non-\ce{CO_2} climate impact are \ce{NO_x}-based ozone production and condensation trails (contrails). Aircraft \ce{NO_x} emissions do not contribute to the greenhouse effect directly, however they act as a precursor to chemical reactions which produce ozone and deplete methane, both of which are potent greenhouse gases. The warming due to tropospheric ozone production tends to outweigh the cooling due to methane reduction, thus resulting in a net warming effect, that constitutes ~16\% of aviation-related global heating \cite{Lee2021}.

In the exceptionally cold and humid conditions experienced in the UTLS region, water vapour can condense and freeze around particulate emissions in the aircraft exhaust to form ice crystals. The trail of visible ice crystals that forms behind the aircraft is known as a contrail. Contrails can be short-lived, dissipating over time with a lifetime of seconds to minutes, or they can persist for many hours, growing with time. The criteria for persistent contrail formation is ambient ice supersaturation. When aircraft generate contrails in ice-supersaturated regions (ISSRs), the high ambient humidity facilitates ice crystal growth, due to uptake of surrounding atmospheric water vapour. Persistent contrails can spread over hundreds of kilometres and transition to cirrus clouds throughout their evolution, resulting in a substantial man-made increase in high ice cloud coverage globally \cite{Karcher2018, Schumann2017}.

The radiative properties of contrail ice crystals are such that they are typically more efficient at absorbing and reflecting infrared radiation emitted from the Earth's surface, than they are at scattering inbound solar radiation, resulting in a global heating effect. The warming induced by aviation contrails and contrail cirrus is on the order of 51\% of total aviation climate impact \cite{Lee2021}. This means that the transient, localised heating effect due to global contrail coverage on a day-to-day basis, is significantly greater than that of all the aircraft \ce{CO_2} emissions that have accumulated since the beginning of civil aviation \cite{Burkhardt2011}.

\subsection{Mitigating Aviation-Induced Climate Change by Modifying Flight Routes}
At cruise altitudes, the degree to which aircraft non-\ce{CO_2} emissions cause warming or cooling is largely determined by the instantaneous atmospheric conditions (i.e. the local chemical and meteorological situation). The atmospheric state is itself, highly variable with respect to time and location, meaning that the same mass of emissions can have very different climate effects, depending on when and where they are released \cite{Matthes2020}. Climate-optimal routing is an operational mitigation concept, proposed to exploit this natural variability in environmental conditions throughout flight. It involves re-routing aircraft to avoid regions of the atmosphere that are particularly sensitive to aircraft emissions. For example, flying around an ISSR region to avoid persistent contrail generation, or flying through airspace in which the chemical composition is less conducive to tropospheric ozone formation, as illustrated in figure \ref{COR}. Modelling efforts have demonstrated the ability to reduce the climate impact of flying by 10--20\%, at the expense of only a few per cent additional fuel burn \cite{Niklass2019a, Matthes2021}.

\begin{figure}[H]
\centering
		\includegraphics[width=0.55\textwidth]{3/COR.pdf}
		%\vspace{-.1cm}
		\caption{Conceptual representation of an aircraft trajectory optimised for minimum climate impact.}
		\label{COR}
\end{figure}

Further to the notion that natural atmospheric fluxes influence aviation's non-\ce{CO_2} climate impact, there is evidence to suggest that artificial perturbations, due to emissions from lingering aircraft exhaust plumes, may also play a significant role \cite{Schlager1997, Jaegle1999}. The overlapping of aircraft plumes can locally elevate chemical mixing ratios of key emission species, resulting in chemical saturation effects which are potentially beneficial from a climate change perspective. Firstly, the ability for \ce{NO_x} emissions to produce ozone become hindered once concentrations surpass a threshold concentration. This means that additional \ce{NO_x} input locally will only serve to reduce the \ce{NO_x}-to-ozone conversion efficiency \cite{Jenkin2002}. Persistent contrail formation is also subject to saturation effects. As a persistent contrail grows, it depletes water vapour from the surrounding atmosphere, thus impeding growth of any subsequent contrails formed in the same region \cite{Unterstrasser2010}. Both these saturation effects lead to a reduction in the resultant non-\ce{CO_2} climate impact of associated flights, compared to the single aircraft scenario. In spite of this, climate-optimal routing studies typically do not consider more than one aircraft, thus neglecting saturation effects caused by residual aircraft emissions. Instead, a range of studies have explored the idea of flying commercial aircraft in formation, as a means of reducing both \ce{CO_2} and non-\ce{CO_2}-related warming. 

Formation flight involves flying one or more aircraft, one behind the other at a nominal separation distance of around 1--2~km. The follower aircraft flies in the upwash of the leader aircraft's wake vortices, which in turn reduces required lift (and hence thrust) for trailing aircraft. As a result, a fuel saving benefit of 5--8\% can be achieved \cite{Xu2014, Bangash2004}. The possibility of exploiting chemical saturation effects through controlled plume overlap has also been explored \cite{Dahlmann2020, Marks2020}. However, no study to date has truly encapsulated the concepts of both climate-optimal routing and formation flight, as two parts of the same optimisation problem. Whilst air traffic data has been used in prior studies of aviation emissions \cite{Filippone2021, Sun2022}, including modelling of contrail formation \cite{Gruber2018}, this is the first paper to the authors' knowledge, that will utilise such datasets to analyse the likelihood of aircraft plume overlap in high-density flight scenarios.